\subsection{Khả năng hội tụ}

\subsubsection{Hồi quy}

\figurename~\ref{fig:regression-loss-curve} mô tả đường cong lỗi của các mô
hình hồi quy sử dụng thuật toán lặp.

\begin{figure}[htbp]
  \centering
  \includegraphics[width=0.9\textwidth]{regression-loss-curve}
  \caption{Biểu đồ đường cong lỗi của các mô hình hồi quy sử dụng
  thuật toán lặp}
  \label{fig:regression-loss-curve}
\end{figure}

Qua biểu đồ, có thể quan sát rằng mô hình MBGD---sử dụng phương pháp lập lịch
tốc độ học dạng step decay---giảm lỗi rất nhanh trong giai đoạn đầu và
hội tụ chỉ sau khoảng 150 epochs. Tuy nhiên, đường cong lỗi của MBGD xuất hiện
các dao động nhẹ quanh xu hướng giảm. Đây là đặc trưng của phương pháp
mini-batch, khi gradient chỉ được ước lượng trên các tập con dữ liệu nên mang
tính nhiễu, dẫn đến quá trình cập nhật không hoàn toàn mượt.

Ngược lại, các mô hình có áp dụng chính quy hoá hội tụ chậm hơn đáng kể, cần
khoảng 1500 epochs để đạt trạng thái ổn định. Một nguyên nhân quan trọng là các
mô hình này không sử dụng các chiến lược lập lịch tốc độ học như
step decay. Do đó, để đảm bảo quá trình học ổn định và tránh dao động
mạnh, tốc độ học phải được chọn nhỏ ngay từ đầu. Điều này làm cho mỗi bước cập
nhật tham số trở nên thận trọng hơn, dẫn đến tốc độ giảm lỗi chậm và kéo dài
thời gian hội tụ.

Bù lại, các mô hình chính quy hoá có đường cong lỗi mượt hơn và ít dao động hơn
so với MBGD, cho thấy quá trình tối ưu ổn định hơn. Nhìn chung, tồn tại sự đánh
đổi rõ rệt giữa tốc độ hội tụ và độ ổn định: MBGD hội tụ nhanh nhưng dao động
nhẹ, trong khi các mô hình chính quy hoá hội tụ chậm hơn nhưng ổn định hơn và
có xu hướng tổng quát hoá tốt hơn.

\subsubsection{Phân lớp}

\figurename~\ref{fig:classification-loss-curve} mô tả đường cong lỗi của các mô
hình hồi quy sử dụng thuật toán lặp.

\begin{figure}[htbp]
  \centering
  \includegraphics[width=0.9\textwidth]{classification-loss-curve}
  \caption{Biểu đồ đường cong lỗi của các mô hình phân lớp sử dụng
  thuật toán lặp}
  \label{fig:classification-loss-curve}
\end{figure}

Từ biểu đồ, có thể thấy rằng cả hai mô hình Softmax và Weighted Softmax đều
hội tụ khá nhanh, chỉ trong khoảng 50 epochs. Điều này phản ánh đặc trưng của
hàm mất mát cross-entropy, với độ dốc lớn ở giai đoạn đầu giúp mô hình nhanh
chóng tiến gần tới vùng cực tiểu.

Tuy nhiên, mô hình Weighted Softmax có đường cong lỗi dao động nhẹ và kém mượt
hơn so với Softmax. Nguyên nhân là do việc áp dụng trọng số lớp làm thay đổi
phân phối gradient giữa các lớp, khiến quá trình cập nhật tham số trở nên kém
ổn định hơn, đặc biệt trong bối cảnh dữ liệu mất cân bằng.

Dù tồn tại dao động, cả hai mô hình đều hội tụ nhanh và không có dấu hiệu phân
kỳ, cho thấy cấu hình huấn luyện là hợp lý. Sự khác biệt về độ mượt của đường
cong lỗi cũng cho thấy một đánh đổi giữa việc xử lý mất cân bằng dữ liệu và độ
ổn định trong quá trình tối ưu.
